\section{Cross-section conventions}\label{app:A.9}

The strength of a dark-matter--electron interaction is quoted as the reference cross-section $\sigmae$ and its momentum dependence as the form factor $\FDM(q)$ of \cref{eq:1.4-defs}, normalized at $q_0 = \alpha m_e = \SI{3.73}{keV}$. For a mediator of mass $m_V$, $\FDM(q) = (m_V^2 + \alpha^2m_e^2)/(m_V^2 + q^2)$, which is 1 for a heavy and $(q_0/q)^2$ for a light mediator~\cite[Eqs.~(11) and (12)]{knapen2021dark}. The form factor multiplies the amplitude, so it enters every rate as $\FDM^2$. Papers define $\sigmae$ and $\FDM(q)$ with different reference points, media and amplitude normalizations.

The momentum $q_0 = \alpha m_e$ is the inverse Bohr radius, typical of bound electrons. A paper that normalizes at another momentum $q_{\mathrm{ref}}$ quotes $\sigma_{\mathrm{ref}} = \sigmae\,\FDM(q_{\mathrm{ref}})^2$; for a light mediator this is $\sigmae\,(q_0/q_{\mathrm{ref}})^4$, so a limit quoted at $q_{\mathrm{ref}} = 2q_0$ is smaller by a factor 16 for the same physics.

Here $\sigmae$ is defined in vacuum: $\tilde V$ is the bare potential between a dark-matter particle and a free electron, and screening, local fields and every other effect of the target are in the response $S$ (\cref{sec:4.4}). The same convention is used in Refs.~\cite[Eq.~(8) and the paragraph after it]{knapen2021dark} and~\cite[below Eq.~(S.8)]{hochberg2021determining}. Some papers put the in-medium dielectric factor into the coupling, as is common for dark photons and for Dirac materials with a background dielectric constant, and quote a cross-section that depends on the target. A target-dependent cross-section of this kind is converted by removing the dielectric factor from the coupling and putting the corresponding screening into the response.

Ref.~\cite[Eq.~(3.8)]{essig2016direct} defines
\begin{equation*}
  \sigmae = \frac{\mu_{\chi e}^2\,|\mathcal{M}_{\mathrm{free}}(q_0)|^2}{16\pi m_\chi^2m_e^2} ,
\end{equation*}
with the relativistically normalized free amplitude $\mathcal{M}_{\mathrm{free}}$, averaged over initial and summed over final spins. Comparison with \cref{eq:1.4-defs} gives the dictionary
\begin{equation}
  |\mathcal{M}_{\mathrm{free}}(q)|^2 = 16\,m_\chi^2m_e^2\,|\tilde V(q)|^2 ,
  \label{eq:A.9-M}
\end{equation}
the factor $4m_\chi m_e$ per amplitude being the ratio of relativistic to nonrelativistic state normalizations. A nonrelativistic effective operator with coefficient $c_1$ in the relativistic normalization therefore has $\tilde V = c_1/(4m_\chi m_e)$ and $\sigmae = \mu_{\chi e}^2|c_1|^2/(16\pi m_\chi^2m_e^2)$.

The spin average is part of this definition. The amplitude squared in \cref{eq:A.9-M} is averaged over the initial spins of the dark matter and the electron and summed over the final spins~\cite[below Eq.~(3.8)]{essig2016direct}. For a spin-independent interaction this average equals the squared amplitude of a single spin state, so $\sigmae$ is the same for scalar and fermionic dark matter with the same $\tilde V$~\cite[below Eq.~(2.2)]{essig2016direct}. The electron spin sum is then not repeated in the target response: there the factor 2 of \cref{app:A.7} counts occupied spin states, not final spins of a free electron~\cite[below Eqs.~(A.28) and (A.29)]{essig2016direct}.

\subsection{Differential cross-section}

The differential cross-section of free elastic scattering in the first Born approximation is \cref{eq:2.2-dsdq},
\begin{equation*}
  \frac{\dd\sigma_T}{\dd q} = \frac{q\,\sigmae}{2\mu_{\chi e}^2v^2}\,\FDM(q)^2, \qquad 0 \le q \le 2\mu_{\chi e}v .
\end{equation*}
For $\FDM = 1$ the differential cross-section integrates to $\sigmae$, and this integral is what the reference cross-section means for a heavy mediator. The same normalization holds for nuclei, $\dd\sigma_N/\dd q = q\sigma_NF_N(q)^2/(2\mu_{\chi N}^2v^2)$ (\cref{eq:2.2-nucleus}). A differential cross-section defined per unit $q^2$ is $\dd\sigma_T/\dd q^2 = (2q)^{-1}\dd\sigma_T/\dd q$; a factor of 2 between two papers often traces to the choice between $q$ and $q^2$.

\subsection{Nucleons and nuclei}

For a mediator coupled to nucleons, the dark-matter--proton cross-section $\sigma_p$ is the analogue of $\sigmae$ with $\mu_{\chi p}$, and the spin-independent nuclear cross-section at small $q$ follows from $\sigma_p$ as $\sigma_N = \sigma_p(\mu_{\chi N}/\mu_{\chi p})^2A^2$ (\cref{eq:1.4-sigmaN}). For unequal proton and neutron couplings, $A^2 \to [Z + (A - Z)y_n/y_p]^2$; Ref.~\cite[Eq.~(104)]{lin2019tasi} writes the equivalent rescaling relative to the neutron coupling. For a dark photon, which couples to protons only, $A^2 \to Z^2$. At larger $q$ the nuclear form factor $F_{\mathrm{Helm}}(q)$ multiplies the amplitude (\cref{sec:2.2}). For a dark photon the electron and proton couplings have equal magnitude $\zeta e$, so for a heavy dark photon, $m_V \gg q_0$, $\sigma_p/\sigmae = (\mu_{\chi p}/\mu_{\chi e})^2$; this ratio is how electron- and nuclear-scattering limits on the same model are compared.

\subsection{Kinetic mixing}

The dark photon mixes with the photon through the term $+(\zeta/2)F_{\mu\nu}F'^{\mu\nu}$ in the Lagrangian density, with $\mathcal{L} \supset -eA_\mu J^\mu_{\mathrm{em}}$ for the photon (\cref{eq:C.9-mix}). The sign of the mixing term follows Ref.~\cite[Eq.~(70)]{lin2019tasi}, which writes $\kappa$ for $\zeta$ and $+eA_\mu J^\mu_{\mathrm{em}}$ for the photon coupling; Ref.~\cite[Eq.~(2.1)]{essig2016direct} writes $\varepsilon$ and mixes with hypercharge, $(\varepsilon/2\cos\theta_W)F_{Y\mu\nu}F'^{\mu\nu}$, where $\theta_W$ is the weak mixing angle. With our signs, the dark photon couples to electrons with $y_e = -\zeta e$ and to protons with $+\zeta e$ (\cref{sec:1.4}). A change of the sign of the mixing term or of the photon coupling flips the sign of $y_e$ and of $\tilde V$, and no rate changes, because every rate contains $\zeta^2$. The symbols used for $\zeta$ in the literature are listed in \cref{app:B.2}.

\subsection{Relativistic dark matter}

The form factor of \cref{eq:1.4-defs} is the nonrelativistic limit, $\omega \ll q$, of the mediator propagator $1/(q^2 - \omega^2 + m_V^2)$ (\cref{app:C.9}). For a light mediator in scattering, $\omega \le qv_{\mathrm{max}}$ and the correction from $\omega^2$ in the propagator is below $10^{-5}$. Dark matter boosted to relativistic speeds is outside this limit: its rates need the full propagator and the coupling to the transverse current, and $\FDM(q)$ alone does not describe them.
